\documentclass[10pt,a4paper,sans]{moderncv}
\moderncvstyle{banking}
\moderncvcolor{blue}
\renewcommand{\sfdefault}{lmss}
\usepackage[a4paper, left=0.6cm, right=0.6cm, top=0.5cm, bottom=0.5cm]{geometry}
\usepackage{enumitem}
\usepackage{hyperref}
\usepackage{fontawesome5}
\usepackage{etoolbox}
\setlist{noitemsep, topsep=0pt, parsep=0pt, partopsep=0pt, leftmargin=*}
\setlength{\parskip}{0pt}
\setlength{\hintscolumnwidth}{2cm}
\renewcommand*{\sectionfont}{\normalsize\bfseries}
\let\oldsection\section
\renewcommand{\section}[1]{\vspace{-1.5ex}\oldsection{#1}\vspace{-0.8ex}}

% Personal Information
\name{Glen}{Muthoka Mutinda}
\phone{+44 7341 625286}
\email{theglenmuthoka@gmail.com}
\social[linkedin]{glenmuthoka}
\social[github]{Bananz0}

\begin{document}

\makecvtitle
\small

\section{Professional Summary}
Final-year BEng Electrical \& Electronics Engineering student at the University of Southampton, with hands-on experience in digital IC physical implementation, FPGA RTL and CPU microarchitecture. Led a TSMC 65nm ring oscillator IC from schematic through custom layout, Mentor Calibre DRC/LVS sign-off and GDSII submission. Seeking a graduate CPU physical implementation role spanning RTL through place-and-route, sign-off and low-power optimisation.

\section{Education}
\cventry[0pt]{2023--2026}{BEng Electrical \& Electronics Engineering, Expected First Class Honours}{University of Southampton}{Southampton, UK}{}{
\textbf{Modules:} Digital Systems Design, Integrated Circuit Design, Embedded Systems, Signal Processing, Control Systems, Network Security. \textbf{Year 2 Project:} TSMC 65nm Ring Oscillator IC (Team Lead). \textbf{Final Year Project:} AI-Driven Optical Authentication using Optical PUFs.\\
\textbf{2022--2023:} Engineering Foundation Programme -- Distinction (ONCAMPUS Global, Southampton).}

\section{Relevant Experience}
\cventry[0pt]{Sep 2023--Jun 2025}{Junior Software Engineer}{ARTEMIS Small Sat-1 Lunar CubeSat Project}{University of Southampton}{}{
\begin{itemize}
\item Developed and debugged C/C++ flight-software components for an 8-engineer student lunar CubeSat across a 2-year mission cycle, operating within MISRA C and NASA-style coding constraints.
\item Implemented real-time telemetry processing under strict timing budgets, contributed on-board sensor-interfacing modules, and built unit-test harnesses to improve integration confidence.
\item Operated daily in Unix-based build and Git version-control environments, reinforcing interface discipline across hardware, software and systems sub-teams.
\end{itemize}}

\section{Selected Hardware \& Implementation Projects}

\cvitem[0pt]{Sep 2023--Jun 2024}{\textbf{TSMC 65nm Real-Time Clock: Ring Oscillator IC} -- Team Lead, 6-engineer team\\
\textit{Cadence S-Edit, T-Spice, L-Edit; Mentor Calibre DRC/LVS}
\begin{itemize}
\item Led a 6-engineer team through the full schematic-to-GDSII physical implementation flow on a TSMC 65nm process, delivering a ring oscillator block for a real-time clock IC.
\item Designed the transistor-level schematic in S-Edit, ran T-Spice simulations for frequency stability and timing closure, and implemented custom L-Edit layout with nanometre-aware floorplanning.
\item Cleared Mentor Calibre DRC and LVS sign-off with zero violations, submitted GDSII for tape-out, and documented the methodology for the following cohort.
\end{itemize}}

\cvitem[0pt]{Apr 2025--Jun 2025}{\textbf{MIPSquare++: 5-Stage Pipelined MIPS Processor Simulator} \hfill \faGithub~\href{https://github.com/Bananz0/MIPSquare}{Bananz0/MIPSquare}\\
\textit{C++23, CMake, Git}
\begin{itemize}
\item Built a cycle-accurate C++23 simulator of a 5-stage MIPS pipeline (Fetch, Decode, Execute, Memory, Write-Back), demonstrating CPU microarchitecture fluency relevant to RTL collaboration.
\item Implemented data forwarding between EX/MEM and MEM/WB stages, load-use hazard detection with automatic stalling, and branch flushing to approach 1 CPI on standard test programs.
\item Engineered a MIPS assembly parser for R-, I- and J-type instructions plus register-file and instruction/data memory models, validated under Unix-based CMake builds.
\end{itemize}}

\cvitem[0pt]{Nov 2024--Dec 2024}{\textbf{16-Stage FIR Notch Filter for Real-Time Audio on FPGA} \hfill \faGithub~\href{https://github.com/Bananz0/16-Stage-FIR-Notch-Filter-for-Real-Time-Audio-Processing-on-FPGA}{Bananz0/FIR-Notch-Filter}\\
\textit{SystemVerilog, Quartus Prime, ModelSim, MATLAB; Altera Cyclone V DE1-SoC}
\begin{itemize}
\item Designed a parameterised 16-stage FIR notch filter in SystemVerilog with a 4-state FSM datapath/controller, achieving 19-cycle latency at 50 MHz on Altera Cyclone V.
\item Generated and validated filter coefficients in MATLAB, then verified RTL functionality in ModelSim before Quartus Prime synthesis and timing-aware FPGA implementation.
\item Synthesised to 2,847 logic elements (approximately 5\% FPGA utilisation) and processed dual-channel 48 kHz stereo audio with -40 dB notch depth.
\end{itemize}}

\cvitem[0pt]{Jan 2025--Mar 2025}{\textbf{WattsApp: Embedded Smart System Management Platform} \hfill \faGithub~\href{https://github.com/Bananz0/WattsApp}{Bananz0/WattsApp}\\
\textit{Embedded C, UART, I2C, ADC, InfluxDB, Grafana; AVR ATMega644p}
\begin{itemize}
\item Built bare-metal embedded C firmware on an AVR ATMega644p, using interrupt-driven sensor polling and UART telemetry under Unix-based toolchains.
\item Implemented I2C and ADC interfaces for live current and voltage acquisition, integrated with an InfluxDB time-series store and Grafana dashboard for real-time visibility.
\item Profiled and optimised the main loop to keep critical-path response within sub-millisecond bounds, demonstrating low-level hardware/software debugging.
\end{itemize}}

\cvitem[0pt]{Sep 2025--Jun 2026}{\textbf{QuantumID: Optical PUF Authentication Framework} -- Lead Developer \& Researcher, Final Year Project\\
\textit{Python, PyTorch, NumPy, SciPy, scikit-learn}
\begin{itemize}
\item Developed an Optical Physical Unclonable Function authentication system in Python using excitation-emission matrix fluorescence spectra and a custom convolutional autoencoder.
\item Engineered a physics-aware preprocessing and augmentation pipeline, then trained reconstruction and contrastive-learning stages to produce latent fingerprints and binary cryptographic hashes.
\item Automated the experiment workflow with NumPy and SciPy under Unix tooling, applying data-driven verification methodology relevant to silicon characterisation flows.
\end{itemize}}

\enlargethispage{3\baselineskip}
\section{Technical Skills}
\cvitem[0pt]{Physical Impl. \& EDA}{TSMC 65nm process exposure, RTL-to-GDSII flow, custom layout, DRC, LVS, place-and-route concepts, timing/area/power trade-offs, nanometre-process awareness, low-power design; Cadence S-Edit/T-Spice/L-Edit, Mentor Calibre, Quartus Prime, ModelSim, Vivado, LTSpice, MATLAB/Simulink.}
\cvitem[0pt]{HDL \& Programming}{SystemVerilog, Verilog, VHDL, RTL/FSM design, pipelined CPU microarchitecture, FPGA synthesis; C, C++, Python, MATLAB, Tcl (basic), Bash, PowerShell, Unix tools, Git, CMake; bare-metal C on AVR/STM32/ESP32, UART/I2C/SPI/ADC.}
\cvitem[0pt]{Memberships \& Status}{IEEE; ISACA; AMRAeS. English/Swahili (Native), Spanish A1. UK Student Visa to 2026; Graduate Route eligible.}

\end{document}
